\chapter{Bone Pain}

\section*{Definition}
Bone pain is a deep, dull, or aching discomfort arising from skeletal tissue, distinct from joint or muscle pain. It may indicate underlying pathology such as fracture, infection, or malignancy when persistent or progressive.

\section*{When to See a Doctor}

\begin{itemize}
 \item Pain is severe, persistent, progressive, or awakens you from sleep
 \item Associated with unexplained fever, night sweats, or weight loss
 \item History of malignancy or pathologic fracture
 \item Bone pain in a child with limping or refusal to bear weight
\end{itemize}

\section*{How It Develops}
Bone pain can arise from injury, inflammation, infection, altered bone metabolism, or a lesion within or around bone. Pain at night or at rest can occur with several conditions and does not by itself diagnose cancer, but persistent or unusual pain should be assessed.

\section*{Causes}
\begin{itemize}
 \item \textbf{Traumatic:} Acute fractures, stress fractures, nonunion
 \item \textbf{Neoplastic:} Primary bone tumors (osteosarcoma, Ewing sarcoma), metastatic disease (breast, prostate, lung, renal)
 \item \textbf{Infectious:} Osteomyelitis or infection near a joint
 \item \textbf{Metabolic:} Osteoporosis with insufficiency fracture, osteomalacia, hyperparathyroidism
 \item \textbf{Hematologic:} Sickle cell vaso-occlusive crisis, leukemia, multiple myeloma
 \item \textbf{Treatment-related:} Some cancer treatments or medicines can cause bone or muscle pain; stopping denosumab can increase fracture risk
\end{itemize}

\section*{Related Symptoms}
\begin{itemize}
 \item Localized tenderness and swelling
 \item Pathologic fracture
 \item Limping or antalgic gait
 \item Erythema and warmth overlying bone
 \item Constitutional symptoms (fever, fatigue, weight loss)
\end{itemize}
\textbf{Important:} Persistent bone pain, swelling, unexplained fever, weight loss, or a fracture without a clear injury warrants prompt medical assessment. Night pain is a warning feature, not proof of malignancy.

\section*{Medical Evaluation}
\begin{itemize}
 \item History and examination guide testing. Plain radiography is often an initial test for focal pain, suspected fracture, or a bone lesion.
 \item MRI or CT is selected for further characterization, occult fracture, infection, nerve involvement, or suspected tumor.
 \item Laboratory evaluation: CBC, ESR, CRP, alkaline phosphatase, calcium, phosphate
 \item Serum protein electrophoresis and urinalysis if multiple myeloma suspected
 \item Bone biopsy for histologic diagnosis if imaging suggests malignancy
 \item Bone scan or PET/CT for staging if metastatic disease is confirmed
\end{itemize}

\section*{Treatment Options}
\begin{itemize}
 \item Pain treatment is individualized and may include acetaminophen/paracetamol or an NSAID when safe; opioids and other adjuvants are reserved for selected cases with medical supervision.
 \item Surgical stabilization of impending or actual pathologic fractures
 \item Radiation therapy for painful osseous metastases
 \item Bisphosphonates or denosumab for osteolytic lesions and osteoporosis
 \item Antibiotic therapy and surgical debridement for osteomyelitis
 \item Treatment of underlying malignancy with chemotherapy, targeted therapy, or immunotherapy
\end{itemize}

\section*{Self-Care and Home Management}
\begin{itemize}
 \item Protect the painful area from activities that worsen symptoms, but avoid prolonged immobilization unless a clinician recommends it.
 \item Use over-the-counter analgesics only as labelled and after checking contraindications; seek advice if pain is persistent or severe.
 \item Cold therapy for acute injuries to reduce inflammation
 \item Avoid weight-bearing until evaluated if fracture suspected
\end{itemize}

\section*{References}
\begin{thebibliography}{99}
\bibitem{bone_pain:ref1} Coleman RE. ``Metastatic bone disease: clinical features, pathophysiology and treatment strategies.'' \textit{Cancer Treat Rev}. 2001;27(3):165-176.
\bibitem{bone_pain:ref2} Sabino MA, Mantyh PW. ``Pathophysiology of bone cancer pain.'' \textit{J Support Oncol}. 2005;3(1):15-24.
\bibitem{bone_pain:ref3} Harrington KD. \textit{Orthopaedic Management of Metastatic Bone Disease}, 2nd ed. Mosby, 2019.
\bibitem{bone_pain:ref4} Lewiecki EM, Binkley N, et al. ``Bone health in cancer survivors.'' \textit{Endocr Rev}. 2022;43(3):495-525.
\bibitem{bone_pain:ref5} National Cancer Institute. ``Primary Bone Cancer.'' Updated 2025. https://www.cancer.gov/types/bone/bone-fact-sheet.
\bibitem{bone_pain:ref6} National Health Service. ``Symptoms of bone cancer.'' Accessed September 2026. https://www.nhs.uk/conditions/bone-cancer/symptoms/.
\end{thebibliography}
